% P1 学术写作版（2026-11-24）：CellCharter研究专题摘要。新文件，不替代任何V1/V2/REV_P0文件。
% 数值来源见 P1_20261124/P1_EVIDENCE_MAP_20261124.md；未编译。
\section*{摘要}
\addcontentsline{toc}{section}{摘要}

空间域识别是空间转录组数据分析的核心任务之一。CellCharter通过多层空间邻域特征聚合与高斯混合模型（GMM）实现空间聚类，并在人类背外侧前额叶皮层（DLPFC）基准数据上表现良好\cite{varrone2024}。然而，对于这类成熟方法，邻域聚合尺度、聚类数确定规则与评价指标三者如何共同影响性能结论，目前尚缺少系统的分析。

本研究以CellCharter原作者提出的方法及其DLPFC基准配置为基础，在3名供体的12张Visium切片上开展系统评价，比较内容包括：0--4层邻域聚合与等维邻域块组合；参考类别数设定、固定K网格与无标签稳定性选K三类聚类数条件；整体分区一致性指标与严格、容差两类空间边界指标。研究还纳入非空间及空间成熟方法作为对照，并在6例乳腺癌患者数据上重复了同一分析流程。

主要结果如下。（1）邻域信息的收益集中在最近邻层：引入第一层邻域后，ARI由0.305升至0.408，且在全部12张切片上方向一致；继续加深只带来小幅增益，并且切片间方向不一致。（2）聚类数设置对结论的影响大于15维之后的尺度差异：固定K时宏平均尺度排序稳定，但逐切片无标签选K后排序改变；在多数七层切片上，ARI在少于参考层数的K处取得最大值。（3）整体一致性与严格边界定位并不同向：随邻域加深，2-hop容差边界F1由0.286升至0.533，而严格边界F1由0.145降至0.118，这种下降主要反映预测边界数量减少带来的召回下降。（4）在乳腺癌队列中，空间邻域特征相对非空间基线在6名患者中方向一致地提高了ARI，但患者间差异很大，与BayesSpace的比较结论也随指标和聚类数规则而变化。

上述结果表明，评价多尺度空间聚类方法时，应同时报告聚类数敏感性、严格与容差边界指标以及供体或患者层面的结果。本研究为成熟空间聚类方法的参数选择与结果解释提供了可复核的实证依据。

\noindent\textbf{关键词：}空间转录组；空间域识别；CellCharter；邻域尺度；聚类数选择；空间边界评价
